\documentclass[11pt]{article}

% review mode (anonymous) for EMNLP 2026 submission
\usepackage[review]{acl}

\usepackage{times}
\usepackage{latexsym}
\usepackage[T1]{fontenc}
\usepackage[utf8]{inputenc}
\usepackage{microtype}
\usepackage{inconsolata}
\usepackage{graphicx}
\usepackage{booktabs}
\usepackage{multirow}
\usepackage{amsmath}
\usepackage{amssymb}

\title{Data-Efficient AI-Code Authorship Detection \\ via Hierarchical Family-Aware Learning}

% Anonymous for EMNLP 2026 review
\author{Anonymous ACL submission}

\begin{document}
\maketitle

%======================================================================
\begin{abstract}
%======================================================================
Detecting which language model authored a code sample is harder than detecting
whether the code is AI-generated at all: state-of-the-art binary detectors
saturate above $99\%$ Macro-F1 on CoDET-M4, while $6$-class authorship sits
at $66\%$. We show that a hierarchical family-aware backbone, when paired with
NTK-target-kernel alignment, attains \textbf{$71.03$ Macro-F1 on CoDET-M4
$6$-class authorship using only $20\%$ of the training data}---a $+4.70$
point improvement over UniXcoder trained on the full set.
The same backbone transfers to DroidCollection 3-class with
\textbf{$89.76$ Weighted-F1 ($+0.98$ over DroidDetectCLS-Large)}.
A family of $14$ architecturally distinct auxiliary losses all cluster above
$70.0$ on the authorship task, suggesting that the gain comes from the shared
hierarchical prior and data-efficient curriculum rather than any single
auxiliary signal. Source-held-out evaluation reveals a persistent
collapse from $0.71$ to $0.36$ on unseen GitHub code, which we characterize
as a source-confounding failure mode.
\end{abstract}

%======================================================================
\section{Introduction}
%======================================================================
\label{sec:intro}

% TODO Day 7: Hook + 3-contribution box.
%   - Hook: AI-code authorship at scale; 6-class is the hard task.
%   - Setup: CoDET-M4 + DroidCollection define the IID/OOD axes.
%   - Failure: 14+ methods cluster around 0.71 IID -> 0.28 OOD-GH.
%   - Our story: hierarchical family prior + NTK alignment, 20% data.
%   - Contributions C1/C2/C3.

%======================================================================
\section{Related Work}
%======================================================================
\label{sec:related}

% TODO Day 8.
%   - AI-code detection: GPTSniffer, CodeGPTSensor, CoDET-M4, DroidCollection,
%     DeTeCtive, CodeMirage.
%   - Authorship as multi-class: hierarchical / contrastive priors.
%   - NTK theory and kernel alignment for OOD generalization.

%======================================================================
\section{Method}
%======================================================================
\label{sec:method}

% TODO Day 3.
\subsection{Problem Formulation}
\label{sec:method-problem}

\subsection{Hierarchical Family Prior}
\label{sec:method-hier}

\subsection{NTK-Aligned Multi-Signal Heads}
\label{sec:method-ntk}

\subsection{Data-Efficient Lean Curriculum}
\label{sec:method-lean}

%======================================================================
\section{Experiments}
%======================================================================
\label{sec:exp}

\subsection{Setup}
\label{sec:exp-setup}

% TODO Day 4: benchmarks, metrics, hardware, training budget.

\subsection{CoDET-M4 Authorship}
\label{sec:exp-codet}

\begin{table}[t]
\centering
\small
\begin{tabular}{llrr}
\toprule
\multicolumn{2}{c}{Method} & Author F1 & $\Delta$ \\
\midrule
\multicolumn{4}{l}{\emph{Paper baselines (full data)}} \\
\multicolumn{2}{l}{UniXcoder} & 66.33 & --- \\
\multicolumn{2}{l}{CodeBERT} & 64.80 & $-1.53$ \\
\multicolumn{2}{l}{CodeT5} & 62.45 & $-3.88$ \\
\midrule
\multicolumn{4}{l}{\emph{Ours, full data}} \\
Exp27 & DeTeCtiveCode & \textbf{71.53} & $+5.20$ \\
\midrule
\multicolumn{4}{l}{\emph{Ours, $20\%$ data (lean)}} \\
Exp\_13 & NTKAlignCode & \textbf{71.03} & $+4.70$ \\
Exp\_06 & FlowCodeDet & 70.90 & $+4.57$ \\
Exp\_09 & EpiplexityCode & 70.78 & $+4.45$ \\
Exp\_12 & AvailPredictivity & 70.35 & $+4.02$ \\
Exp\_07 & SAMFlatCode & 70.22 & $+3.89$ \\
\bottomrule
\end{tabular}
\caption{CoDET-M4 6-class authorship Macro-F1. Twenty methods cluster
above $70.0$ in our family; UniXcoder is the strongest paper baseline at
$66.33$. Lean methods use only $20\%$ of training data.}
\label{tab:codet-author}
\end{table}

\subsection{DroidCollection T3 Cross-Bench}
\label{sec:exp-droid}

\begin{table}[t]
\centering
\small
\begin{tabular}{llr}
\toprule
\multicolumn{2}{c}{Method} & T3 W-F1 \\
\midrule
\multicolumn{3}{l}{\emph{Paper baselines (full data)}} \\
\multicolumn{2}{l}{DroidDetectCLS-Large} & 88.78 \\
\multicolumn{2}{l}{DroidDetectCLS-Base} & 86.76 \\
\multicolumn{2}{l}{Fast-DetectGPT (zero-shot)} & 64.54 \\
\midrule
\multicolumn{3}{l}{\emph{Ours, $20\%$ data (lean)}} \\
Exp\_04 & PoincareGenealogy & \textbf{89.76} \\
Exp\_06 & FlowCodeDet & 89.34 \\
Exp\_09 & EpiplexityCode & 89.26 \\
Exp\_03 & TokenStatRAG & 88.94 \\
\bottomrule
\end{tabular}
\caption{DroidCollection 3-class Weighted-F1 on the full test split.
Multiple methods exceed the full-data DroidDetectCLS-Large baseline
using only $20\%$ training data.}
\label{tab:droid-t3}
\end{table}

\subsection{Data Efficiency}
\label{sec:exp-efficiency}

% TODO Day 5: data-efficiency curve (5/10/20/50/100% training data).

\subsection{Ablation}
\label{sec:exp-ablation}

% TODO Day 5-6: ablation matrix from each Exp_N file.

%======================================================================
\section{Analysis}
%======================================================================
\label{sec:analysis}

% TODO Day 9.
\subsection{Source-Confounding Failure Mode}
\label{sec:analysis-source}

% Train -> test gap 0.71 -> 0.36 on held-out GitHub source.

\subsection{Sibling Generator Confusion}
\label{sec:analysis-siblings}

% Qwen <-> Nxcode pair is the primary IID bottleneck.

\subsection{Cross-Method Robustness}
\label{sec:analysis-robustness}

% Author F1 vs OOD-gh F1 scatter across 14+ methods (Figure 2).

%======================================================================
\section{Limitations}
%======================================================================
\label{sec:limitations}

% TODO Day 10. OOD-gh still hard; AICD-Bench out of scope.

%======================================================================
\section{Conclusion}
%======================================================================
\label{sec:conclusion}

% TODO Day 10.

% Bibliography
\bibliography{custom}

\end{document}
